\section{Removing the enhancement of a flag}\label{sec:descent}

We continue over $\mathbb C$.  Write $N_B=R_u(B)$, $T=B/N_B$, and
\[
 q:\cA^+=\Bun_{G,N_B,x}(X)\longrightarrow
 \cA=\Bun_{G,B,x}(X).
\]
This is a $T$-torsor: an enhancement of a fixed Borel reduction is a
trivialization of its induced $T$-torsor.  Set $r=\dim T$.  Our objective
is to show that a nonzero eigencomplex on $\cA^+$ produces one on $\cA$.
The direct image $q_!$ need not preserve nonvanishing for arbitrary
monodromy along its fibers.  We will prove that the unipotent boundary
monodromy of the eigenvalue forces unipotent monodromy along those
fibers, which is sufficient.

The central sheaves of Gaitsgory and their universal monodromic version
of Dhillon--Taylor provide the local identity that relates these two
monodromies.  The main point of this section is to globalize that
identity with its tensor and individual-factor monodromy maps.  All
universal local systems used in this argument are Betti sheaves.  The
final direct image and perverse truncation are performed on the original
geometric adic objects.

\subsection{Monodromy along a torsor and the universal unit}

We call a complex on a $T$-torsor \emph{monodromic along the torsor} if,
locally on its base, it is constructible for a product stratification
$\{S_\alpha\times T\}$ and its restriction to a contractible patch in
the $T$-factor is pulled back from a slice.  This permits arbitrary
monodromy in $\pi_1(T)=X_*(T)$; it does not require an equivariant
structure.  The condition is preserved by base change on the base.

\begin{lemma}\label{lem:descent-monodromic}
Let $Q$ be a locally bounded constructible complex on $\cA^+$ with
singular support in the generic nilpotent cone.  Its Betti realization
is monodromic along $q$.
\end{lemma}

\begin{proof}
The cotangent description in Section~\ref{sec:geometry} identifies the
nilpotent cone on $\cA^+$ with the smooth pullback of the cone on $\cA$:
a nilpotent residue lying in a Borel Lie algebra has zero toral
component.  In particular every covector in this cone annihilates the
tangent to a $T$-fiber.

Take a smooth finite-type chart $D\to\cA$ and trivialize the pulled-back
torsor, after further localization on $D$.  The singular support on
$D\times T$ lies in $C_D\times 0_T$, where $C_D\subset T^*D$ is the
ordinary-level cone of Proposition~\ref{prop:cone-dimension}.  Choose a
microlocal stratification of $D$ whose conormals contain $C_D$.
The singular-support criterion for constructibility then applies to
the product stratification $\{S_\alpha\times T\}$; this is precisely
the criterion used in the proof of \cite[Proposition~6.3.2]{NY}.
On a contractible patch of $T$, restriction to a slice is an
equivalence for sheaves constructible with respect to this product
stratification.  Consequently the complex, including its extension
data between strata, is pulled back from the slice there.
The descriptions agree on overlaps by locally constant continuation.
They establish the asserted condition on the torsor and after every
further base change.  In particular $X_*(T)$ acts by continuation on
the stalk cohomology of the restriction to each fiber.
\end{proof}

Fix orientations of the compact real torus in $T(\mathbb C)$ and the
corresponding positive cocharacter loops.  Put
\[
 \Gamma=X_*(T),\qquad
 R_T=E[\Gamma]=\cO(\widehat T).
\]
Let $\mathscr L_T$ be the local system on $T$ with fiber $R_T$ and
regular monodromy.  Thus its pullback to the universal cover of the
compact torus is constant, with fiber the finitely supported functions
on the deck group.  It is a weakly constructible Betti sheaf, with
infinite-dimensional stalks.  We use ordinary direct image in the
convolution of such kernels, and its unit is
\begin{equation}\label{eq:descent-unit}
 \mathbf 1_T=\mathscr L_T[r].
\end{equation}
For a fixed output enhancement $e$, write a varying input enhancement
as $e d^{-1}$.  The coordinate $d$ is the enhancement difference in our
convolution convention.  With this choice the endomorphisms $R_T$ of
the unit act by the monodromy of the input, without inversion.

\begin{lemma}\label{lem:descent-unit}
Convolution with \eqref{eq:descent-unit} acts as the identity on
complexes monodromic along a $T$-torsor.  This identification is natural
under base change, compatible with successive convolutions, and
identifies the $R_T$-action with enhancement monodromy.
\end{lemma}

\begin{proof}
First take a local system with fiber $W$, and write $M_1,\ldots,M_r$
for its commuting monodromy operators in a basis of $\Gamma$.
For fixed output, the integrand has monodromies
$z_i\otimes M_i^{-1}$ on $R_T\otimes_E W$.  Ordinary cohomology of
$T$ is computed by the cohomological Koszul complex
\[
 K^\bullet\bigl(z_1M_1^{-1}-1,\ldots,z_rM_r^{-1}-1;
                   R_T\otimes_E W\bigr).
\]
Untwisting the diagonal $\Gamma$-action identifies this with the
Koszul complex for $z_1-1,\ldots,z_r-1$ on the free module with
underlying fiber $W$.  Its only cohomology is $W$ in degree $r$.
The shift in \eqref{eq:descent-unit} therefore makes the integral $W$
in degree zero, and the residual $z_i$ acts by $M_i$.

There is a useful geometric description of this identification.
Remove the contractible real radial directions of $T$.  The free local
system is the direct image with finite supports from the universal
cover $\mathbb R^r$ of the compact torus.  Integration over that
compact torus becomes compactly supported integration over
$\mathbb R^r$.  On the cover the other factor is pulled back from its
fiber by continuation from the zero lift.  The orientation identifies
$R\Gamma_c(\mathbb R^r,E)[r]$ with $E$.
This construction also works for complexes and relatively over a
stratified base, so it proves naturality and base change.

For two enhancement differences $d_1,d_2$, convolution replaces them
by their product $d_2d_1$.  On their real covers this is addition of
lifted angles.  The change from the two lifted differences to the
first lifted difference and the lifted total difference preserves
product orientation.  Integrating in this order agrees with the two
successive integrations; the continuation of the input follows the
sum of the same two paths.  This proves compatibility with the unit
convolution isomorphism, including its shifts and $R_T$-action.
\end{proof}

\subsection{The local central-sheaf identity}

Let $LG$ denote loops in a formal curve coordinate, and let
$\mathrm{Iw}$ and $\mathrm{Iw}^0$ be the inverse images of $B$ and
$N_B$ under evaluation $L^+G\to G$.  We use the Betti universal
monodromic Hecke category on $LG/\mathrm{Iw}^0$, with the
pro-unipotent equivariance convention of \cite{DT}.  Its unit is
\eqref{eq:descent-unit} on the identity stratum $T$.
The local result we need is the monoidal central-sheaf construction
\[
 Z:\Rep(\Gd)\longrightarrow\mathcal H^{\mathrm{mon}},
 \qquad V\longmapsto Z(V),
\]
together with its tensor automorphism $m_V$ given by nearby-cycle
monodromy.  These are constructed in \cite[Section~6, especially
Proposition~6.3.1]{DT}, using the degeneration introduced in
\cite{GaitsgoryCentral}.

We record explicitly the consequence of the local monodromy formula
that will be globalized.  If $\chi_V$ is the character of $V$ restricted
to $\widehat T$, then
\begin{equation}\label{eq:descent-local-trace}
 \left(
 \mathbf 1_T\xrightarrow{\mathrm{coev}}
 Z(V)\star Z(V^*)
 \xrightarrow{m_V\star\mathrm{id}}
 Z(V)\star Z(V^*)
 \xrightarrow{\mathrm{ev}}\mathbf 1_T
 \right)=\chi_V\quad\text{in }\operatorname{End}(\mathbf 1_T)=R_T.
\end{equation}
Here evaluation in the indicated order is the image of evaluation in
the symmetric representation category.  To deduce the equation from
\cite[Proposition~9.3.1(a)--(c)]{DT}, apply the faithful monoidal
Wakimoto associated-graded functor of \cite[Proposition~5.3.1]{DT}.
On the weight summand indexed by
$\nu\in X^*(\widehat T)=\Gamma$, the monodromy is multiplication by
the torus function $e^\nu$.  The trace on the associated graded is
therefore $\sum_\nu(\dim V_\nu)e^\nu=\chi_V$.  The associated-graded
functor is faithful on endomorphisms of the unit: on its sole identity
stratum it is the regular $R_T$-module and retains the multiplication
endomorphism.  Thus equality of these graded endomorphisms proves
\eqref{eq:descent-local-trace}.  Opposite conventions invert both the
torus and the boundary loop; we keep the convention fixed above.

\subsection{Specialization through a torus torsor}

To use the local trace formula globally, we must commute nearby cycles
with a direct image whose nonproper fibers are enhancement tori.
The following observation is the required substitute for properness.

\begin{lemma}\label{lem:descent-radial}
Let $\Delta$ be an analytic disc, let $Y\to\Delta$ be a stratified
analytic space, and factor a morphism over $\Delta$ as
\[
 P\xrightarrow{\pi} C\xrightarrow{b}Y,
\]
where $\pi$ is a $T$-torsor on the entire family and $b$ is proper.
Work locally on finite-dimensional, paracompact charts.  Let $F$ on
$P|_{\Delta^*}$ be weakly constructible and, in torsor charts,
constructible for product stratifications with the $T$-factor.
Then the exchange map
\begin{equation}\label{eq:descent-push-exchange}
 \Psi(b\pi)_*F\longrightarrow(b_0\pi_0)_*\Psi F
\end{equation}
is an isomorphism.  The same reduction gives the base-change and
projection-formula isomorphisms for these direct images when the
additional factors are pulled back from their targets.  These
identifications respect composition.
\end{lemma}

\begin{proof}
Let $T_c$ be the compact real form of $T$, and let
$A\simeq\mathbb R^r$ be its positive real radial subgroup.  A
reduction of the structure group from $T$ to $T_c$ exists on the
paracompact charts.  Equivalently, quotient the torsor by $A$:
\[
 P\xrightarrow{\rho}\overline P=P/A\xrightarrow{\bar\pi}C.
\]
Here $\overline P$ is a topological quotient, and we use topological
nearby cycles, defined by the universal cover of the punctured disc.
The first map has contractible real fibers, and $\bar\pi$ is a
locally trivial compact-torus bundle, hence proper.  These are
factorizations over the full disc, including its center.

The hypothesis implies that $F=\rho^*\overline F$ for a sheaf
$\overline F$ on $\overline P|_{\Delta^*}$.  This may be checked in a
torsor chart, where it is the restriction equivalence for the product
stratification and a contractible radial factor.  Such local
descriptions glue because that restriction equivalence is fully
faithful.  In a trivialization extending across zero, nearby cycles
are computed in the product of a radial patch with a neighborhood in
$\overline P$.  Contracting the radial patch gives
\[
 \Psi F=\rho_0^*\Psi\overline F,
 \qquad R\rho_*\rho^*\overline F=\overline F.
\]
These identities hold also after base change and after tensoring by
a factor from the target.  Now $b\bar\pi$ is proper, so ordinary
proper base change gives \eqref{eq:descent-push-exchange}.
It also proves the other asserted formulas after the radial factors
have been removed.  Every identification is the natural exchange
map, as can be seen using restriction and direct image from the
universal cover of $\Delta^*$.  They therefore compose for composite
maps.  The same argument works for products of enhancement tori.
\end{proof}

\subsection{The central action and its individual monodromies}

We now compare the local central action with the global eigenvalue.  A choice of
lift to the universal cover of a punctured disc at $x$ defines a
nearby fiber $V_{\sigma,\partial}$ of every $V_\sigma$, compatibly with
tensors; let $M_{V,\partial}$ be its positive-loop monodromy.

\begin{lemma}[Global action of the central kernels]
\label{lem:central-globalization}
Let $Q$ be a locally bounded constructible Betti complex on $\cA^+$,
monodromic along $q$, equipped with coherent Hecke eigenisomorphisms
on $U$ with eigenvalue $\sigma$.  The kernels $Z(V)$ act at $x$ on
$Q$, by ordinary direct image in bounded affine Hecke
correspondences, and there are isomorphisms
\begin{equation}\label{eq:descent-central-eigen}
 Z(V)\star_x Q\simeq Q\otimes_E V_{\sigma,\partial}.
\end{equation}
They are natural in $V$, compatible with the tensor unit and
successive convolutions, and carry $m_V\star_x\mathrm{id}_Q$ to
$\mathrm{id}_Q\otimes M_{V,\partial}$.  In a convolution of several
central kernels this comparison holds for the monodromy on each
individual factor.  Consequently the action on $Q$ of
\eqref{eq:descent-local-trace} is
\begin{equation}\label{eq:descent-global-trace}
 \chi_V(\text{enhancement monodromy})
 =\operatorname{Tr}(M_{V,\partial})\,\mathrm{id}_Q.
\end{equation}
\end{lemma}

\begin{proof}
The trace in \eqref{eq:descent-local-trace} applies monodromy to the
$Z(V)$ factor alone.  We must therefore construct the central
eigenisomorphism together with its tensor compatibility, retaining the
monodromy on each factor.  We first specialize one moving action, then
compare successive actions on a common correspondence, and finally
identify the individual monodromies and evaluate the trace.

\smallskip\noindent
\emph{The moving action and its specialization.}
Choose a coordinate disc $\Delta$ around $x$, with $x=0$.  Consider
the correspondence $\mathfrak H_\Delta$ whose points consist of a
leg $z\in\Delta$, two enhanced bundles $(P_0,\eta_0)$ and
$(P_1,\eta_1)$, and an isomorphism
\[
 P_0|_{X\setminus\{z\}}\simeq P_1|_{X\setminus\{z\}}.
\]
The two enhancements at $x$ are independent, including when
$z\ne0$.  Write $p$ for the input map and $o$ for the output-and-leg
map.  We use bounded versions of this correspondence throughout.

In a frame of $P_0$ that takes $\eta_0$ to the standard enhanced
flag, the local family is the Gaitsgory family of \cite[Section~6]{DT}.
Its fiber for $z\ne0$ is
\[
 \Gr_G\times G/N_B,
\]
and its fiber at zero is $LG/\mathrm{Iw}^0$.
For $V\in\Rep(\Gd)$, let $K_V$ on the punctured family be the Satake
kernel on the Grassmannian factor, externally tensored with
$\mathbf 1_T$ on the locally closed subspace
$T=B/N_B\subset G/N_B$, extended by zero.  Thus its support requires
the two ordinary flags to agree under the modification, while its
universal local system records their enhancement difference.
By the local construction,
\begin{equation}\label{eq:descent-kernel-specialization}
 \Psi K_V=Z(V).
\end{equation}
We use unsuspended geometric nearby cycles on relative kernels; the
moving-curve perverse shift and its inverse in perverse nearby cycles
have both been removed.  This is the normalization in which the
moving tensor-unit Hecke functor is exterior product with $E_\Delta$.

These descriptions descend from the frames to the global
correspondence.  Indeed changes of input frame are regular loop
changes whose value at $x$ lies in $N_B$.  Off zero the Satake factor
has its positive-loop equivariance and the flag factor has the
corresponding $N_B$-invariance.  The descent isomorphisms and their
cocycle identities specialize by smooth pullback.  At zero their
group is $\mathrm{Iw}^0$, which is exactly the required equivariance
for \eqref{eq:descent-kernel-specialization}.
On a fixed bound all this can be carried out with finite jets: use a
sufficiently thick neighborhood of the sum of the graphs of $x$ and
$z$, so the same frame space works at collision.  These frame spaces
are smooth.  Off zero, descent of morphisms uses the specified
Satake and flag equivariance.  At collision the jet quotients of
$\mathrm{Iw}^0$ are unipotent, hence contractible in Betti topology;
restriction along their group nerves is fully faithful.  Thus the
descent retains the maps of kernels, not only their isomorphism
classes.  Gluing bundles off the disc
identifies these local descriptions with the global ones.

Over $\Delta^*$ form the action
\[
 A_V(Q)=o_*\bigl(Q\,\widetilde\boxtimes K_V\bigr),
\]
where the twisted product includes the fixed relative normalization.
On the support of $K_V$, ordinary flags agree.  Hence the
correspondence is the usual enhancement-preserving Hecke
correspondence together with a varying input enhancement.  Its
$T$-integration is the unit calculation of
Lemma~\ref{lem:descent-unit}.  It gives, with its functorial unit
identification,
\begin{equation}\label{eq:descent-punctured-action}
 A_V(Q)\simeq \Hecke_V(Q)|_{\Delta^*}
       \simeq Q\boxtimes V_\sigma|_{\Delta^*}.
\end{equation}

We next justify specialization of the direct image in this equation.
The obstruction to properness is entirely the input enhancement.
Forgetting that enhancement, while retaining its ordinary flag,
is a $T$-torsor over the \emph{whole} bounded correspondence over
$\Delta$.  Its quotient is proper over output and leg: after fixing
the output bundle and enhancement, a bounded input modification is
in a proper Grassmannian bound, and the ordinary input flag is a
point of the proper flag bundle over it.  This description uses full
flags over the bound and remains valid at $z=0$.

Both factors of the integrand are relatively monodromic in this
input-enhancement torsor.  For $Q$ this is the base-change-stable
condition in Lemma~\ref{lem:descent-monodromic}, or the hypothesis of
the present lemma.  For $K_V$ over $\Delta^*$, its ordinary-flag
condition and Grassmannian factor are independent of the enhancement;
its remaining factor is a local system in the enhancement difference.
Thus their product satisfies Lemma~\ref{lem:descent-radial}.
It follows that nearby cycles commute with $o_*$.  This argument
uses one radial quotient extending across zero, so it also proves
the asserted comparison at collision, where properness of $o$
itself is unavailable.

On the correspondence, nearby cycles of the integrand are its
twisted product with $Z(V)$.  To check this assertion, pass to the
input frames just described.  There the input is pulled back from
the bundle factor and the kernel from the local degeneration.
The nearby-cycle exterior-product map is an isomorphism.  More
generally the same statement holds when the input is a family $D$
over $\Delta^*$ that is $Q$ tensored with a lisse complex there:
in the untwisted notation of a frame chart, it is the natural map
\begin{equation}\label{eq:descent-tensor-exchange}
 p_0^*\Psi D\otimes\Psi K_V
       \longrightarrow \Psi(p^*D\otimes K_V).
\end{equation}
Trivialize the lisse factor on the universal cover of $\Delta^*$ to
reduce to the same exterior-product calculation.

For precision concerning coefficients, these are calculations in
weakly constructible Betti sheaves on finite-dimensional bounds.
Adapt a complex stratification to the kernels and input.  The local
Milnor data have finite triangulations, and the exterior comparison
is computed by cellular cochains on nearby fibers in product
neighborhoods, using the same lift to the universal cover of the
puncture.  Each complex has finitely many cells; the usual K\"unneth
map therefore works for the vector-space stalks of the universal
local system as well.  There is no inverse limit of finite-rank adic
systems in this assertion.  The computation is natural in maps of
the factors and respects their monodromy; smooth frame changes
preserve it.  This is also the local calculation used in
\cite[Proposition~6.3.1]{DT}.

Write $A_V$ for the punctured-family action above and
\[
 B_V(R)=Z(V)\star_x R
\]
for its collision action.  For an input family $D$ to which
\eqref{eq:descent-tensor-exchange} applies, let
\begin{equation}\label{eq:descent-c-map}
 c_V(D):B_V(\Psi D)\longrightarrow\Psi A_V(D)
\end{equation}
be the tensor-exchange map followed by the inverse of the push
exchange.  In frame notation it is the composite
\[
 o_{0,*}(p_0^*\Psi D\otimes Z(V))
 \longrightarrow o_{0,*}\Psi(p^*D\otimes K_V)
 \xrightarrow{\sim}\Psi o_*(p^*D\otimes K_V).
\]
The twists in this formula are the same on both sides.  The preceding
calculations prove that $c_V(D)$ is an isomorphism for $D$ constant
with value $Q$, and for $D=Q\boxtimes L$ with $L$ a finite-rank local
system on $\Delta^*$.

For constant input, compose $c_V(Q)$ with nearby cycles of
\eqref{eq:descent-punctured-action}.  This gives
\eqref{eq:descent-central-eigen} and identifies $m_V$ with
$M_{V,\partial}$, since $Q$ is constant in the disc direction.
We now compare these maps for tensor products.  Compatibility with
total nearby-cycle monodromy would only identify the product of the
factor monodromies; the trace requires each factor separately.

\smallskip\noindent
\emph{Successive actions and the tensor comparison.}
For a second representation $W$, the eigenisomorphism identifies
$A_W(Q)$ with $Q\boxtimes W_\sigma|_{\Delta^*}$, so it is among the
inputs for which $c_V$ is an isomorphism.  We first record how this
comparison treats an extra lisse factor.  For $D=Q\boxtimes L$, set
$L_\partial=\Psi L$.  Projection formula
and the exterior comparison give the following commutative square:
\begin{equation}\label{eq:descent-lisse-square}
\begin{CD}
 B_V\bigl(\Psi(Q\boxtimes L)\bigr)
   @>{c_V(Q\boxtimes L)}>> \Psi A_V(Q\boxtimes L)\\
 @V{\sim}VV @VV{\sim}V\\
 B_V(Q)\otimes L_\partial
   @>{c_V(Q)\otimes\mathrm{id}}>>
       \Psi A_V(Q)\otimes L_\partial.
\end{CD}
\end{equation}
Indeed, on the universal cover of $\Delta^*$ the factor $L$ is
constant.  There both routes are the tensor-exchange map for $Q$
and $K_V$ tensored with its fiber, followed by the same push exchange.
Finite rank permits that fiber to pass through all the direct images.
The construction is equivariant for deck transformations, so the
square descends, and it is natural in every local-system endomorphism
of $L$.  In particular it respects its boundary monodromy as an
endomorphism on this factor alone: that monodromy commutes with the
cyclic fundamental group of $\Delta^*$ and hence is an endomorphism
of $L$, not merely an automorphism of its nearby fiber.

It remains to identify the resulting successive comparison with the
one for the convolution of kernels.  Let
$\mathfrak H^{(2)}_{W,V}$ parametrize two bounded modifications at
$z\in\Delta$,
\[
 (P_0,\eta_0)\longrightarrow(P_1,\eta_1)
                  \longrightarrow(P_2,\eta_2),
\]
with each enhancement independent.  The first bound supports $K_W$
and the second supports $K_V$.  There are two forgetting maps.
The map $f:\mathfrak H^{(2)}_{W,V}\to\mathfrak H_V$ forgets
$(P_0,\eta_0)$ and the first modification; it is the base change of
the output map for $\mathfrak H_W$.  Forgetting only $\eta_0$, while
retaining its ordinary flag $\beta_0$, factors $f$ as a $T$-torsor
followed by a proper map, over the whole disc.  For its projection
formula the second kernel is pulled back from $\mathfrak H_V$.

The other map
$g:\mathfrak H^{(2)}_{W,V}\to\mathfrak H_{V\otimes W}$ forgets
$(P_1,\eta_1)$ and composes the two modifications; take a large enough
bound on the target to contain their composites.  First forget only
$\eta_1$, retaining $P_1$ and its ordinary flag $\beta_1$.  This is
again a $T$-torsor on the entire family.  The remaining map is proper.
For fixed endpoints and composite modification, the possible bounded
$P_1$ form the closed convolution fiber inside a proper Grassmannian
bound: the two relative-position bounds are closed conditions.
Choosing $\beta_1$ adds a proper $G/B$-bundle.  This description
is valid also at $z=0$, where the composite is an ordinary local
modification with the endpoint enhancements retained.

The integrand $Q\,\widetilde\boxtimes K_W
\,\widetilde\boxtimes K_V$ is monodromic in the torus forgotten by
either map.  For $f$ this is the already proved input-enhancement
calculation, with the second kernel pulled back from the target.
For $g$ the factor $Q$ is independent of $\eta_1$.  Over $\Delta^*$,
changing $\eta_1$ changes the two enhancement differences inversely;
the two universal local systems therefore remain a local system in
that variable.  Their ordinary-flag conditions and Grassmannian
factors do not depend on $\eta_1$.  This proves relative
monodromicity in the actual intermediate-enhancement torsor, in
frames, without requiring an equivariant structure on $Q$.
Lemma~\ref{lem:descent-radial} now applies to $f$ and $g$, their base
changes, and the required projection formulas.  The same uniform
radial quotient applies to their specialized integrands by the
nearby-cycle tensor comparison.

For $z\ne0$, the two unit integrations give a convolution isomorphism
\[
 b:A_VA_W(Q)\xrightarrow{\sim}A_{V\otimes W}(Q),
\]
where tensor factors are written in action order.  Forgetting the
intermediate bundle first gives exactly this map: on the two torus
covers, replace the two lifted enhancement differences by the first
lifted difference and their lifted total difference.  The coordinate
change preserves product orientation.  Integration in the first
variable is the free-unit convolution map of
Lemma~\ref{lem:descent-unit}, and its continuation agrees with the
successive unit integrations.

At zero the DT convolution comparison gives
\[
 b_0:B_VB_W(Q)\xrightarrow{\sim}B_{V\otimes W}(Q).
\]
It is the action of the specific monoidal kernel map
$Z(V)\star Z(W)\simeq Z(V\otimes W)$ of
\cite[Section~6.3]{DT}: in the successive frames,
\cite[Lemma~6.3.2]{DT} identifies the nearby cycles of the twisted
product of the two kernels with their specialized twisted product,
and pushing by $g$ is precisely their construction of this map.
Exterior comparison with $Q$ and frame descent preserve that map.

The equality of the two ways to compare actions is the square
\begin{equation}\label{eq:descent-convolution-square}
\begin{CD}
 B_VB_W(Q) @>{b_0}>> B_{V\otimes W}(Q)\\
 @V{c^{\mathrm{seq}}_{V,W}(Q)}VV
     @VV{c_{V\otimes W}(Q)}V\\
 \Psi\bigl(A_VA_W(Q)\bigr)
    @>{\Psi(b)}>> \Psi A_{V\otimes W}(Q),
\end{CD}
\end{equation}
where
\[
 c^{\mathrm{seq}}_{V,W}(Q)
   =c_V(A_W(Q))\circ B_V(c_W(Q)).
\]
All domains in this expression are identified by $\Psi Q=Q$ for
the constant input family.  To prove the square, on
$\mathfrak H^{(2)}_{W,V}$ take the natural map from the product of
the pulled-back nearby cycles of $Q,K_W,K_V$ to the nearby cycles
of their twisted product.  It is an isomorphism by the DT
comparison just cited and the exterior comparison with $Q$.
Push this map first through $f$ or first through $g$.  Through $f$,
base change and projection formula identify it with
$c_V(A_W(Q))\circ B_V(c_W(Q))$.  Through $g$, they identify it with
the action of the DT monoidal map followed by
$c_{V\otimes W}(Q)$.  These identifications give the same map after
the final output push: pullback-and-tensor exchange composes for
iterated tensors, push exchange composes for composite maps, and
their intervening projection-formula square commutes.  Equivalently
these are the exchange squares for restriction and direct image
from the punctured-disc cover.  The uniform torus factorizations
above justify every exchange, so this verifies
\eqref{eq:descent-convolution-square} also for the nonproper maps.

\smallskip\noindent
\emph{Individual monodromy and the trace.}
Apply this square using the eigenisomorphism
$A_W(Q)\simeq Q\boxtimes W_\sigma|_{\Delta^*}$.
The second input is covered by \eqref{eq:descent-lisse-square}:
it is still monodromic along the enhancement, and its extra factor
is lisse on $\Delta^*$.  The single-step map $c_W(Q)$ identifies
$m_W$ with $M_{W,\partial}$.  Applying $B_V$ to that identification
and then \eqref{eq:descent-lisse-square} carries this endomorphism
to the $W_{\sigma,\partial}$ factor alone.  This is legitimate
independently of simultaneous disc monodromy, since the latter
square is natural in the local-system endomorphism
$M_{W,\partial}$ of $W_\sigma|_{\Delta^*}$.
For the outer factor $m_V$, the bottom row of
\eqref{eq:descent-lisse-square} is $c_V(Q)$ tensored with the
identity of $W_{\sigma,\partial}$, so its single-step identification
gives $M_{V,\partial}\otimes\mathrm{id}$.
Thus under the action comparison the two individual endomorphisms
are respectively $\mathrm{id}\otimes M_{W,\partial}$ and
$M_{V,\partial}\otimes\mathrm{id}$.

The convolution square and the original tensor compatibility of
the eigenisomorphisms show that these successive identifications
are the one for $V\otimes W$.  Naturality of all maps of kernels
proves compatibility with representation morphisms, including
coevaluation and evaluation.  The unit compatibility is exactly
Lemma~\ref{lem:descent-unit}.  Iterating the same square gives any
number of factors, and permutation and fusion maps are the
corresponding Satake kernel maps transported by these natural
exchanges.

We may therefore evaluate \eqref{eq:descent-local-trace} on $Q$.
Its coevaluation and evaluation become the ordinary insertion and
contraction for $V_{\sigma,\partial}$, and the middle map becomes
$M_{V,\partial}$ on its indicated factor.  The resulting scalar is
$\operatorname{Tr}(M_{V,\partial})$.  On the other hand, the right
side of \eqref{eq:descent-local-trace} acts through the $R_T$-action
identified in Lemma~\ref{lem:descent-unit}.  This proves
\eqref{eq:descent-global-trace} and the lemma.
\end{proof}

\subsection{Unipotence and nonvanishing of the direct image}

\begin{proposition}\label{prop:torus-descent}
Let $\sigma$ be a Zariski-dense geometric adic parameter on $U$
whose local monodromy at $x$ is unipotent.  If $\cA^+$ carries a
nonzero locally bounded constructible geometric adic Hecke
eigencomplex for $\sigma$, with the coherent tensor and fusion data,
then $\cA$ carries a nonzero locally constructible perverse geometric
adic Hecke eigensheaf for the same parameter and with the same
coherences.  Its singular support lies in the parabolic global
nilpotent cone.
\end{proposition}

\begin{proof}
By Proposition~\ref{prop:perverse-eigen}, a nonzero perverse cohomology
of the given eigencomplex is again an eigenobject.  Denote it by
$Q$.  Its singular support is nilpotent by
Theorem~\ref{thm:nilpotent-detection}; hence its Betti realization
satisfies Lemma~\ref{lem:descent-monodromic}.
Lemma~\ref{lem:central-globalization} gives, for every algebraic
representation $V$ of $\Gd$,
\begin{equation}\label{eq:descent-character-identity}
 \chi_V(\text{enhancement monodromy})=(\dim V)\,\mathrm{id}_Q,
\end{equation}
because the boundary monodromy of $V_\sigma$ is unipotent.

Restrict to any torsor fiber, and to any finite-dimensional stalk
cohomology space of $Q$ on it.  The commuting enhancement
monodromies have joint generalized eigencharacters
$a:\Gamma\to E^\times$, that is, points $a\in\widehat T(E)$.
Equation~\eqref{eq:descent-character-identity} says
$\chi_V(a)=\chi_V(1)$ for every $V$.  In characteristic zero the
representation characters span $\cO(\widehat T)^W$ and separate
semisimple conjugacy classes.  Thus $a$ and $1$ have the same image
in the finite Weyl quotient $\widehat T/W$.  Its fiber through $1$
is the singleton $\{1\}$, so $a=1$.
Every enhancement monodromy operator is therefore unipotent on
every stalk cohomology space.  This argument requires no uniform
bound on Jordan blocks as the point of the bundle stack varies.

Return to geometric adic sheaves and set $F=q_!Q$.
The morphism $q$ has finite type and fixed relative dimension, so
$F$ is locally bounded constructible.  We prove $F\ne0$ by one
fiber.  Choose a point of $\cA^+$ with nonzero stalk, and let $a$ be
its image in $\cA$.  On the fiber $T=q^{-1}(a)$ the restriction
$Q_T$ has locally constant finite-dimensional cohomology sheaves.
Choose the largest $j$ for which $\mathcal H^j(Q_T)\ne0$, and
write $W$ for its fiber.  This is a nonzero finite-dimensional
representation of $\Gamma$ with commuting unipotent monodromies.
Its coinvariants $W_\Gamma$ are nonzero.  Indeed the augmentation
ideal of $E[\Gamma]$ acts nilpotently on $W$: its finitely many
commuting nilpotent generators have a common finite nilpotence
bound on this particular space.  If $W_\Gamma$ vanished, successive
powers of that ideal would give $W=0$.

Poincar\'e duality on the torus identifies
\[
 H_c^{2r}(T,\mathcal H^j(Q_T))\simeq W_\Gamma(-r)\ne0.
\]
In the compact-support hypercohomology spectral sequence, the term
of bidegree $(2r,j)$ has no incoming differential, since all higher
cohomology sheaves vanish, and no outgoing differential, since
compactly supported cohomology vanishes above degree $2r$.
It yields a nonzero class in $H_c^{2r+j}(T,Q_T)$.
Comparison with Betti sheaves and base change for $q_!$ now show
$F_a\ne0$.

Finally, away from $x$ the Hecke correspondence transports both the
ordinary flag and its enhancement.  Thus the enhanced
correspondence is the pullback of the ordinary one along $q$, and
the bounded output map is proper.  Base change and the projection
formula give natural isomorphisms
\[
 \Hecke_{I,(V_i)}(q_!Q)
 \simeq(q\times\mathrm{id}_{U^I})_!
                   \Hecke_{I,(V_i)}(Q)
 \simeq q_!Q\boxtimes\mathop{\boxtimes}_{i\in I}(V_i)_\sigma.
\]
The same correspondence calculation works on the successive
modification spaces and on every collision diagonal in $U^I$.
The exchange maps are natural, so they preserve the unit, tensor,
permutation, and fusion diagrams.  Hence $F$ is a nonzero coherent
Hecke eigencomplex on $\cA$.  A nonzero perverse cohomology of $F$
is an eigensheaf by Proposition~\ref{prop:perverse-eigen}, and its
singular support is in $\Lambda_{\mathrm{par}}$ by
Theorem~\ref{thm:nilpotent-detection}.
\end{proof}

The descent argument uses unipotence of boundary monodromy, without
a regularity requirement.  In the application to Theorem~\ref{thm:main},
the prescribed regular-unipotent boundary satisfies this hypothesis.
